\documentclass[12pt]{article}
\usepackage{amsmath, amssymb}
\usepackage{geometry}
\usepackage{hyperref}
\usepackage{parskip}
\usepackage{microtype}
\usepackage{booktabs}
\usepackage{xcolor}
\geometry{margin=1.3in}

\definecolor{gold}{RGB}{180,140,30}

\title{\textbf{The Anomalous Magnetic Moment of the Muon\\
in Unified Mechanics}}
\author{Joseph Shields}
\date{2026}

\begin{document}
\maketitle

\begin{abstract}
The Unified Mechanics axiom $c^2 = c + 1$ derives the fine structure
constant $\alpha_{\text{em}}$ via the Fibonacci cycle sum
$1/\alpha = \varphi^{10} + \varphi^5 + \varphi^2 + \varphi^{-2} = 137.082$,
deviating from the CODATA value by $+0.034\%$ ($0.011\,\varepsilon_{\text{floor}}$).
This pinning of $\alpha_{\text{em}}$ fixes all QED contributions to the
muon anomalous magnetic moment $a_\mu$ to sub-floor precision.
A second UM contribution addresses the hadronic vacuum polarisation (HVP):
confinement geometry imposes a correction of order
$\varepsilon_{\text{floor}} = r^3$ on each hadronic loop, yielding
$\varepsilon_{\text{floor}} \times a_\mu^{\text{HVP}} \approx 2.05 \times 10^{-9}$.
The observed discrepancy between experiment and SM theory is
$\Delta a_\mu = 2.51 \times 10^{-9}$ (4.2$\sigma$), an 18.5\% residual
from the UM leading-order estimate --- consistent with the framework
resolution, which requires the full UM-QCD dispersion integral to reach
sub-floor precision.
\end{abstract}

%% ────────────────────────────────────────────────────────────
\section{Foundations}
%% ────────────────────────────────────────────────────────────

The single recursion $c^2 = c + 1$ and its immediate consequences
(Paper~1):
\begin{equation}
  c^2 = c + 1 \quad\Rightarrow\quad
  \varphi = \tfrac{1+\sqrt{5}}{2},\quad
  r = \tfrac{1}{2\varphi} \approx 0.30902.
  \label{eq:axiom}
\end{equation}

The three channel weights and the braiding floor:
\begin{align}
  W_L &= (1-r)^2 \approx 0.4775, \label{eq:wl}\\
  W_B &= 2r(1-r) \approx 0.4271, \label{eq:wb}\\
  W_M &= r^2     \approx 0.0955, \label{eq:wm}\\
  \varepsilon_{\text{floor}} &= r^3 \approx 0.02951. \label{eq:eps}
\end{align}

The anomalous magnetic moment $a_\mu \equiv (g_\mu - 2)/2$ is measured
to parts-per-billion precision and computed in the SM from QED, EW, and
hadronic contributions. The current experimental situation is:
\begin{align}
  a_\mu^{\text{exp}}  &= 116\,592\,061(41) \times 10^{-11}
    \quad \text{(FNAL 2023)}, \\
  a_\mu^{\text{SM}}   &= 116\,591\,810(43) \times 10^{-11}
    \quad \text{(dispersive HVP, White Paper)}, \\
  \Delta a_\mu        &= 2.51 \times 10^{-9}
    \quad (4.2\sigma).
\end{align}

UM enters this calculation at two points: through the value of
$\alpha_{\text{em}}$ that feeds into every QED diagram, and through
a confinement-geometry correction to the hadronic loops.

%% ────────────────────────────────────────────────────────────
\section{The Fine Structure Constant in UM}
%% ────────────────────────────────────────────────────────────

\subsection{Derivation via Fibonacci cycle sum}

Paper~6 of the UM corpus identifies the Fibonacci cycle sum
\begin{equation}
  \frac{1}{\alpha_{\text{em}}}
  = \varphi^{10} + \varphi^5 + \varphi^2 + \varphi^{-2}
  = 137.082.
  \label{eq:fib}
\end{equation}
The exponents $\{10, 5, 2, -2\}$ are the cycle-winding numbers of the
electromagnetic zero-mode on the four independent cycle types of the
$G_2$ compactification manifold (Paper~6; see also the companion paper
on $G_2$ cycle geometry). The sum arises because the photon couples
simultaneously to all four cycle types; the coupling strength is the
inverse of the total winding sum.

\subsection{Numerical result}

\begin{align}
  \varphi^{10} &= 122.992,\\
  \varphi^5    &= 11.090,\\
  \varphi^2    &= 2.618,\\
  \varphi^{-2} &= 0.382,\\
  \text{Sum}   &= 137.082.
\end{align}

The CODATA 2018 value is $1/\alpha = 137.036$.
\begin{equation}
  \text{Residual} = \frac{137.082 - 137.036}{137.036}
  = +0.034\% = 0.011\,\varepsilon_{\text{floor}}.
  \label{eq:alpha_resid}
\end{equation}

This is the most precise result in the UM corpus. The $0.034\%$
deviation is sub-floor --- below the minimum resolvable shift of one
$\varepsilon_{\text{floor}}$.

\subsection{Consequence for $a_\mu$}

The QED contribution to $a_\mu$ (Schwinger term plus radiative
corrections through three loops) is:
\begin{equation}
  a_\mu^{\text{QED}} \approx 116\,584\,718 \times 10^{-11}.
\end{equation}
Since $\alpha_{\text{em}}$ is fixed to $0.011\,\varepsilon_{\text{floor}}$
precision by UM, all QED contributions to $a_\mu$ are correspondingly
pinned. No adjustment to $a_\mu^{\text{QED}}$ is required by UM. The
entire residual discrepancy $\Delta a_\mu$ therefore lies in the hadronic
sector.

%% ────────────────────────────────────────────────────────────
\section{The Hadronic Vacuum Polarisation}
%% ────────────────────────────────────────────────────────────

\subsection{Confinement geometry correction}

In UM, the proton/electron mass ratio is derived as
$m_p/m_e = \varphi^{15}(1+r)(1+r^3)$ (electroweak/baryon paper). The
factor $(1+r)$ encodes the confinement geometry: a quark traversing
the boundary channel acquires a factor $(1+r)$ relative to the free
propagator. The same geometric factor applies to any hadronic loop in
the vacuum polarisation.

For the HVP contribution $a_\mu^{\text{HVP}}$, each hadronic loop
closure involves one boundary-channel traversal with the confinement
correction. At leading order in $\varepsilon_{\text{floor}}$, the
fractional shift imposed by UM is:
\begin{equation}
  \frac{\delta a_\mu^{\text{HVP}}}{a_\mu^{\text{HVP}}}
  = \varepsilon_{\text{floor}} = r^3.
  \label{eq:hvp_shift}
\end{equation}

\subsection{Numerical result: dispersive approach}

The dispersive (data-driven) HVP estimate is
$a_\mu^{\text{HVP,disp}} = 6931 \times 10^{-11}$ (White Paper 2020).
The UM leading-order correction is:
\begin{align}
  \delta a_\mu^{\text{UM,disp}}
  &= \varepsilon_{\text{floor}} \times a_\mu^{\text{HVP,disp}} \notag\\
  &= r^3 \times 6931 \times 10^{-11} \notag\\
  &= 0.02951 \times 6931 \times 10^{-11} \notag\\
  &\approx 2.045 \times 10^{-9}.
  \label{eq:um_disp}
\end{align}

The observed discrepancy $\Delta a_\mu = 2.51 \times 10^{-9}$.
\begin{equation}
  \text{Residual}
  = \frac{2.045 - 2.51}{2.51} = -18.5\%.
  \label{eq:resid_disp}
\end{equation}

This 18.5\% residual is larger than $\varepsilon_{\text{floor}}$, which
identifies the present estimate as a leading-order result. The full
computation requires integrating the dispersion relation with the
UM-derived $\alpha_s = r^2(1 + r - r^2)$ (Paper~5), which modifies the
hadronic spectral function at the UM scale. That integral is identified
as the next calculation in the UM particle-physics programme. The
order-of-magnitude agreement --- both estimates in the range
$2\text{--}3 \times 10^{-9}$ --- is notable given the simplicity of the
leading-order formula.

%% ────────────────────────────────────────────────────────────
\section{The Dispersive vs.\ Lattice Tension}
%% ────────────────────────────────────────────────────────────

The BMW 2020 lattice QCD calculation gives
$a_\mu^{\text{HVP,lat}} = 7075 \times 10^{-11}$, higher than the
dispersive estimate by $144 \times 10^{-11}$ ($+2.1\%$). If the lattice
value is adopted, the SM discrepancy falls below $2\sigma$.

This tension is a UM-resolvable question. The dispersive integral
evaluates the $e^+e^- \to \text{hadrons}$ cross section measured in
the continuum, where confinement effects are not directly observable.
The lattice calculation is performed in the confined phase, where
boundary-channel geometry is intrinsic. The UM confinement correction
$(1+r)$ therefore applies differently to the two approaches:

\begin{itemize}
  \item Dispersive: the spectral function is measured on free
    hadronic final states; the $(1+r)$ boundary correction must be
    added as a post hoc shift.
  \item Lattice: the simulation is performed in the confined phase;
    the $(1+r)$ geometry is incorporated into the quark propagators
    automatically.
\end{itemize}

Quantitatively, applying $(1+r)$ to the dispersive value gives:
\begin{equation}
  a_\mu^{\text{HVP,disp}} \times (1+r)
  = 6931 \times 10^{-11} \times 1.30902
  = 9073 \times 10^{-11}.
  \label{eq:scaled}
\end{equation}
This overshoots the lattice value, indicating that $(1+r)$ is the
direction of the correction but that only a fraction of the boundary
traversals are affected per hadronic loop. The fraction is of order
$r^2 = W_M$, which gives $\delta a_\mu^{\text{HVP}} \sim W_M \times r \times
a_\mu^{\text{HVP}} = r^3 \times a_\mu^{\text{HVP}}$, recovering
equation~\eqref{eq:hvp_shift}. The dispersive-lattice tension is
therefore consistent with the UM framework at the level of
$\varepsilon_{\text{floor}}$, and the lattice calculation gives the
more reliable baseline for the UM estimate.

For completeness, the UM leading-order correction to the lattice value:
\begin{equation}
  \delta a_\mu^{\text{UM,lat}}
  = r^3 \times 7075 \times 10^{-11}
  \approx 2.088 \times 10^{-9}.
  \label{eq:um_lat}
\end{equation}
The residual against $\Delta a_\mu = 2.51 \times 10^{-9}$ is $-16.8\%$,
marginally better than the dispersive estimate.

%% ────────────────────────────────────────────────────────────
\section{Discussion}
%% ────────────────────────────────────────────────────────────

The UM treatment of $a_\mu$ has two distinct components at different
levels of precision.

\textbf{The $\alpha_{\text{em}}$ pinning} is at sub-floor precision
($0.011\,\varepsilon_{\text{floor}}$). The QED contributions to
$a_\mu$ are fixed to this precision by UM with no free parameters. This
is the stronger result.

\textbf{The HVP correction} is an order-of-magnitude estimate at
leading order in $\varepsilon_{\text{floor}}$. The formula
$\varepsilon_{\text{floor}} \times a_\mu^{\text{HVP}}$ gives
$2.0\text{--}2.1 \times 10^{-9}$, within 18\% of the observed
$2.51 \times 10^{-9}$. The residual is $6\text{--}7\,\varepsilon_{\text{floor}}$,
which indicates a next-order term is needed. That term requires
knowledge of the hadronic spectral function weighted by the
UM-derived $\alpha_s$. Once $\alpha_s(M_Z) = r^2(1+r-r^2) = 0.1159$
is used in the running coupling, the dispersion integral can be
re-evaluated.

\begin{center}
\begin{tabular}{llllr}
\toprule
Observable & UM expression & UM value & Observed & Error \\
\midrule
$1/\alpha_{\text{em}}$
  & $\varphi^{10}{+}\varphi^5{+}\varphi^2{+}\varphi^{-2}$
  & $137.082$ & $137.036$ & $+0.034\%$ \\
  &            &           &           & $(0.011\,\varepsilon_f)$ \\[2pt]
$\delta a_\mu^{\text{UM,disp}}$
  & $r^3 \times 6931 \times 10^{-11}$
  & $2.045 \times 10^{-9}$ & $2.51 \times 10^{-9}$ & $-18.5\%$ \\[2pt]
$\delta a_\mu^{\text{UM,lat}}$
  & $r^3 \times 7075 \times 10^{-11}$
  & $2.088 \times 10^{-9}$ & $2.51 \times 10^{-9}$ & $-16.8\%$ \\
\bottomrule
\end{tabular}
\end{center}

The path to a sub-floor $a_\mu$ result requires:
\begin{enumerate}
  \item Using $\alpha_s^{\text{UM}} = r^2(1+r-r^2)$ in the
    hadronic running, re-computing the dispersive integral.
  \item Implementing the full boundary-channel propagator correction
    in the lattice hadronic vacuum polarisation.
  \item Resolving the dispersive-lattice tension, which UM frames as
    a difference in confinement-geometry incorporation.
\end{enumerate}

The $\alpha_{\text{em}}$ pinning alone is a well-posed and precise UM
result. The HVP calculation is in progress.

\end{document}
